% TOP_PROBLEM: {"id": 3327, "title": "Consecutive non-orientable embedding obstructions", "category_id": 3, "category": {"id": 3, "name": "graph_theory", "display_name": "Graph Theory", "description": "Problems involving graphs, networks, and their properties.", "slug": "graph-theory", "order_index": 3, "created_at": "2026-07-31T15:26:25.671Z"}, "status": "open", "problem_number": "OPG-172", "set_id": 12, "statusQualification": "Quantifies genus g>=1 and explains why the title’s consecutive version is equivalent to existence for any two distinct nonorientable surfaces."}
% FIRST_POSTED: 2026-10-01
% ENTRY_KIND: statement_only
% UnsolvedMath ID 3327; source catalogue code OPG-172
% !TeX program = lualatex
% Definition and sources only; this file is not an LLM solution attempt.
\documentclass[11pt,a4paper]{article}
\usepackage[margin=25mm]{geometry}
\usepackage{fontspec}
\setmainfont{lmroman10-regular.otf}[BoldFont=lmroman10-bold.otf,
 ItalicFont=lmroman10-italic.otf,BoldItalicFont=lmroman10-bolditalic.otf]
\usepackage{amsmath,amssymb,mathtools}
\usepackage{unicode-math}
\setmathfont{latinmodern-math.otf}
\AtBeginDocument{\let\setminus\smallsetminus}
\usepackage{xurl,hyperref}
\hypersetup{colorlinks=true,urlcolor=blue}
\setcounter{secnumdepth}{0}
\setlength{\parindent}{0pt}
\setlength{\parskip}{5pt}
\setlength{\emergencystretch}{3em}
\begin{document}
\section*{Consecutive non-orientable embedding obstructions}\label{3327}
{\small UnsolvedMath ID 3327 · OPG-172 · Graph Theory · OpenGarden}
\subsection{Definitions and mathematical statement}
For an integer $g\ge1$, let $N_g$ be the connected closed nonorientable surface that is the connected sum of $g$ projective planes. A finite graph embeds in $N_g$ if its vertices can be placed at distinct points and its edges drawn as arcs meeting only at common endpoints. Embeddings are not required to be cellular.

A minor of a finite simple graph is obtained by deleting vertices or edges and contracting edges, subsequently discarding loops and merging parallel edges. A proper minor is one not isomorphic to the original graph. Call $G$ a minor-minimal obstruction for $N_g$ when $G$ does not embed in $N_g$ but every proper minor of $G$ does.

The question asks whether there exist a finite graph $G$ and an integer $g\ge1$ such that $G$ is a minor-minimal obstruction for both $N_g$ and $N_{g+1}$. Equivalently, one requires
\[
 G\not\hookrightarrow N_{g+1},\qquad
 H\hookrightarrow N_g\quad\text{for every proper minor }H\text{ of }G.
\]
Embeddability in $N_g$ implies embeddability in $N_{g+1}$, so these conditions give both obstruction properties.

This consecutive-surface form also captures the wording asking for any two distinct nonorientable surfaces. If a graph is an obstruction for $N_i$ and $N_j$ with $i<j$, monotonicity makes it an obstruction for each intervening surface, and in particular for $N_i,N_{i+1}$. Thus no particular genus is fixed in advance, and the existence quantifier over $g$ is essential.
\subsection{Short English statement}
Can a single finite graph be a forbidden minor for two consecutive nonorientable surfaces?
\subsection{Sources}
\begin{itemize}
\item[S1] ULAM.AI and the UnsolvedMath Contributors, supplied problems.json, record 3327 (OPG-172), OpenGarden. Catalogue identity and source transcription; CC BY 4.0.
\url{https://huggingface.co/datasets/ulamai/UnsolvedMath}.
\item[S2] Open Problem Garden, Consecutive non-orientable embedding obstructions. Original problem and discussion.
\url{http://www.openproblemgarden.org/op/consecutive_non_orientable_embedding_obstructions}.
\end{itemize}
\end{document}
